\textbf{Phần 1. Trắc nghiệm lựa chọn}
\setcounterpageref{ex}{0}

\begin{ex} %[5D7V36]
Một lớp học có $35$ học sinh, trong đó có $21$ bạn nữ. Tỉ số phần trăm của số bạn nam so với số học sinh cả lớp là
\choice
{$30\%$}
{\True $40\%$}
{$50\%$}
{$60\%$}
\loigiai{
Số bạn nam trong lớp là:
$35 - 21 = 14\text{ (bạn)}$.\\
Tỉ số phần trăm của số bạn nam so với cả lớp là:
$14 : 35 = 0{,}4 = 40\%$.\\
\textbf{Đáp án của câu hỏi là B}
}
\end{ex}

\begin{ex} %[5D7V36]
Viết tỉ số $\dfrac{18}{25}$ dưới dạng tỉ số phần trăm ta được
\choice
{$18\%$}
{$36\%$}
{\True $72\%$}
{$90\%$}
\loigiai{
Ta có: $\dfrac{18}{25} = \dfrac{18 \times 4}{25 \times 4} = \dfrac{72}{100} = 72\%$.\\
\textbf{Đáp án của câu hỏi là C}
}
\end{ex}

\begin{ex} %[5D7V36]
Một khu vườn có $120$ cây cam và $80$ cây bưởi. Tỉ số phần trăm của số cây bưởi so với tổng số cây trong khu vườn là
\choice
{$30\%$}
{\True $40\%$}
{$60\%$}
{$66{,}6\%$}
\loigiai{
Tổng số cây trong vườn là:
$120 + 80 = 200\text{ (cây)}$.\\
Tỉ số phần trăm của số cây bưởi so với tổng số cây là:
$80 : 200 = 0{,}4 = 40\%$.\\
\textbf{Đáp án của câu hỏi là B}
}
\end{ex}

\begin{ex} %[5D7V36]
Một chiếc áo có giá niêm yết là $200\,000$ đồng. Nhân dịp lễ, cửa hàng giảm giá $15\%$. Giá của chiếc áo sau khi giảm giá là
\choice
{$150\,000\text{ đồng}$}
{\True $170\,000\text{ đồng}$}
{$180\,000\text{ đồng}$}
{$185\,000\text{ đồng}$}
\loigiai{
Số tiền được giảm giá là:
$200\,000 \times 15\% = 30\,000\text{ (đồng)}$.\\
Giá của chiếc áo sau khi giảm giá là:
$200\,000 - 30\,000 = 170\,000\text{ (đồng)}$.\\
\textbf{Đáp án của câu hỏi là B}
}
\end{ex}

\begin{ex} %[5D7V36]
Số học sinh xuất sắc của một trường tiểu học là $180$ em, chiếm $24\%$ tổng số học sinh toàn trường. Số học sinh của toàn trường đó là
\choice
{$600\text{ học sinh}$}
{$720\text{ học sinh}$}
{\True $750\text{ học sinh}$}
{$800\text{ học sinh}$}
\loigiai{
Số học sinh của toàn trường là:
$180 : 24\% = 180 : 0{,}24 = 750\text{ (học sinh)}$.\\
\textbf{Đáp án của câu hỏi là C}
}
\end{ex}

\begin{ex} %[5D7V36]
Khi tăng một số thêm $25\%$ rồi lại giảm số mới nhận được đi $20\%$ thì số cuối cùng so với số ban đầu
\choice
{Tăng thêm $5\%$}
{Giảm đi $5\%$}
{\True Không thay đổi}
{Giảm đi $1\%$}
\loigiai{
Gọi số ban đầu là $A$.\\
Khi tăng $25\%$, số mới là: $A \times (1 + 0{,}25) = 1{,}25 \times A$.\\
Khi giảm số mới đi $20\%$, số cuối cùng là:
$(1{,}25 \times A) \times (1 - 0{,}2) = 1{,}25 \times 0{,}8 \times A = 1 \times A = A$.\\
Vậy số cuối cùng không thay đổi so với số ban đầu.\\
\textbf{Đáp án của câu hỏi là C}
}
\end{ex}

\begin{ex} %[5D7V36]
Tỉ số giữa số $A$ và số $B$ là $\dfrac{3}{5}$. Tỉ số phần trăm của số $A$ so với số $B$ là
\choice
{$30\%$}
{$50\%$}
{\True $60\%$}
{$166{,}67\%$}
\loigiai{
Tỉ số phần trăm của số $A$ so với số $B$ là:
$3 : 5 = 0{,}6 = 60\%$.\\
\textbf{Đáp án của câu hỏi là C}
}
\end{ex}

\begin{ex} %[5D7V36]
Lượng nước trong hạt cà phê tươi chiếm $20\%$ khối lượng. Bác An đem phơi $500\text{ kg}$ cà phê tươi để lấy hạt khô nguyên chất (coi như nước bay hơi hoàn toàn). Khối lượng hạt cà phê khô thu được là
\choice
{$350\text{ kg}$}
{\True $400\text{ kg}$}
{$420\text{ kg}$}
{$450\text{ kg}$}
\loigiai{
Hạt cà phê thuần túy chiếm số phần trăm là:
$100\% - 20\% = 80\%$.\\
Khối lượng hạt cà phê khô thu được là:
$500 \times 80\% = 400\text{ (kg)}$.\\
\textbf{Đáp án của câu hỏi là B}
}
\end{ex}

\begin{ex} %[5D7V36]
Một mảnh đất hình chữ nhật có chiều dài $25\text{ m}$, chiều rộng bằng $60\%$ chiều dài. Diện tích của mảnh đất đó là
\choice
{$300\text{ m}^2$}
{\True $375\text{ m}^2$}
{$400\text{ m}^2$}
{$450\text{ m}^2$}
\loigiai{
Chiều rộng mảnh đất là:
$25 \times 60\% = 15\text{ (m)}$.\\
Diện tích mảnh đất là:
$25 \times 15 = 375\text{ (m}^2\text{)}$.\\
\textbf{Đáp án của câu hỏi là B}
}
\end{ex}

\begin{ex} %[5D7V36]
Tỉ số giữa số học sinh giỏi và số học sinh khá của một trường là $\dfrac{4}{7}$. Biết số học sinh giỏi ít hơn số học sinh khá là $45$ bạn. Số học sinh giỏi của trường đó là
\choice
{$50\text{ bạn}$}
{\True $60\text{ bạn}$}
{$70\text{ bạn}$}
{$105\text{ bạn}$}
\loigiai{
Hiệu số phần bằng nhau là:
$7 - 4 = 3\text{ (phần)}$.\\
Số học sinh giỏi là:
$45 : 3 \times 4 = 60\text{ (bạn)}$.\\
\textbf{Đáp án của câu hỏi là B}
}
\end{ex}

\begin{ex} %[5D7V36]
Bác Ba gửi tiết kiệm $50\,000\,000$ đồng vào ngân hàng với lãi suất $6\%$ một năm. Sau một năm, tổng số tiền cả gốc và lãi bác Ba nhận được là
\choice
{$51\,500\,000\text{ đồng}$}
{\True $53\,000\,000\text{ đồng}$}
{$56\,000\,000\text{ đồng}$}
{$54\,500\,000\text{ đồng}$}
\loigiai{
Số tiền lãi sau một năm là:
$50\,000\,000 \times 6\% = 3\,000\,000\text{ (đồng)}$.\\
Tổng số tiền nhận được là:
$50\,000\,000 + 3\,000\,000 = 53\,000\,000\text{ (đồng)}$.\\
\textbf{Đáp án của câu hỏi là B}
}
\end{ex}

\begin{ex} %[5D7V36]
Tỉ số phần trăm của hai số $12$ và $16$ là
\choice
{$60\%$}
{$70\%$}
{\True $75\%$}
{$80\%$}
\loigiai{
Ta có: $12 : 16 = 0{,}75 = 75\%$.\\
\textbf{Đáp án của câu hỏi là C}
}
\end{ex}

\textbf{Phần 2. Trắc nghiệm đúng sai}
\setcounterpageref{ex}{0}

\begin{ex} %[5D7V36]
Một trường tiểu học có $800$ học sinh. Trong đó số học sinh nữ chiếm $52\%$ tổng số học sinh toàn trường.
\choiceTF
{\True Tỉ số phần trăm của số học sinh nam so với học sinh toàn trường là $48\%$}
{\True Trường đó có $416$ học sinh nữ}
{\True Trường đó có $384$ học sinh nam}
{Số học sinh nữ nhiều hơn số học sinh nam là $40$ học sinh}
\loigiai{
\begin{itemchoice}
\item Đúng. Tỉ số phần trăm nam: $100\% - 52\% = 48\%$.
\item Đúng. Số học sinh nữ: $800 \times 52\% = 416\text{ (học sinh)}$.
\item Đúng. Số học sinh nam: $800 - 416 = 384\text{ (học sinh)}$.
\item Sai. Nữ nhiều hơn nam là: $416 - 384 = 32\text{ (học sinh)} 
eq 40\text{ học sinh}$.
\end{itemchoice}
}
\end{ex}

\begin{ex} %[5D7V36]
Một cửa hàng mua một chiếc xe đạp với giá vốn $1\,500\,000$ đồng. Cửa hàng bán lại với mức lãi $20\%$ so với giá vốn.
\choiceTF
{\True Số tiền lãi khi bán chiếc xe đạp là $300\,000$ đồng}
{\True Giá bán của chiếc xe đạp là $1\,800\,000$ đồng}
{\True Tiền lãi bằng $\dfrac{1}{5}$ tiền vốn ban đầu}
{Nếu muốn thu được mức lãi $25\%$ giá vốn thì cửa hàng phải bán xe với giá $1\,900\,000$ đồng}
\loigiai{
\begin{itemchoice}
\item Đúng. Tiền lãi: $1\,500\,000 \times 20\% = 300\,000\text{ (đồng)}$.
\item Đúng. Giá bán: $1\,500\,000 + 300\,000 = 1\,800\,000\text{ (đồng)}$.
\item Đúng. Vì $20\% = \dfrac{20}{100} = \dfrac{1}{5}$.
\item Sai. Muốn lãi $25\%$ giá vốn thì giá bán phải là: $1\,500\,000 \times (1 + 25\%) = 1\,875\,000\text{ (đồng)} 
eq 1\,900\,000\text{ đồng}$.
\end{itemchoice}
}
\end{ex}

\textbf{Phần 3. Trắc nghiệm trả lời ngắn}
\setcounterpageref{ex}{0}

\begin{ex} %[5D7V36]
Lớp 5A có $30$ bạn tham gia thi đấu cờ vua, trong đó có $18$ bạn nam. Tính tỉ số phần trăm của số bạn nữ so với tổng số học sinh của cả đội.
\shortans{40}
\loigiai{
Số bạn nữ là:
$30 - 18 = 12\text{ (bạn)}$.\\
Tỉ số phần trăm của số bạn nữ là:
$12 : 30 = 0{,}4 = 40\%$.
}
\end{ex}

\begin{ex} %[5D7V36]
Tìm một số biết rằng $30\%$ của số đó có giá trị bằng $75$.
\shortans{250}
\loigiai{
Số cần tìm là:
$75 : 30\% = 75 : 0{,}3 = 250$.
}
\end{ex}

\begin{ex} %[5D7V36]
Một cuộn dây dài $80\text{ m}$. Người ta cắt đi $35\%$ chiều dài của cuộn dây đó. Chiều dài đoạn dây còn lại là bao nhiêu mét?
\shortans{52}
\loigiai{
Tỉ số phần trăm của đoạn dây còn lại là:
$100\% - 35\% = 65\%$.\\
Chiều dài đoạn dây còn lại là:
$80 \times 65\% = 52\text{ (m)}$.
}
\end{ex}

\begin{ex} %[5D7V36]
Một đàn gà có $150$ con gồm gà trống và gà mái. Biết rằng số gà trống bằng $\dfrac{2}{3}$ số gà mái. Số gà mái trong đàn là bao nhiêu con?
\shortans{90}
\loigiai{
Tổng số phần bằng nhau là:
$2 + 3 = 5\text{ (phần)}$.\\
Số gà mái là:
$150 : 5 \times 3 = 90\text{ (con)}$.
}
\end{ex}

\begin{ex} %[5D7V36]
Một mảnh đất hình chữ nhật có chu vi $70\text{ m}$. Chiều rộng bằng $\dfrac{3}{4}$ chiều dài. Diện tích của mảnh đất đó là bao nhiêu mét vuông?
\shortans{300}
\loigiai{
Nửa chu vi mảnh đất là:
$70 : 2 = 35\text{ (m)}$.\\
Chiều rộng mảnh đất là:
$35 : (3 + 4) \times 3 = 15\text{ (m)}$.\\
Chiều dài mảnh đất là:
$35 - 15 = 20\text{ (m)}$.\\
Diện tích mảnh đất là:
$20 \times 15 = 300\text{ (m}^2\text{)}$.
}
\end{ex}

\begin{ex} %[5D7V36]
Một chiếc quạt điện có giá niêm yết ban đầu là $500\,000$ đồng. Nhân dịp khuyến mãi, cửa hàng giảm giá $12\%$. Hỏi người mua được giảm giá bao nhiêu nghìn đồng?
\shortans{60}
\loigiai{
Số tiền được giảm giá là:
$500\,000 \times 12\% = 60\,000\text{ (đồng)}$.\\
Đổi: $60\,000\text{ đồng} = 60\text{ nghìn đồng}$.
}
\end{ex}
